\FloatBarrier
\section{FT outer flange}

The seventh beamline study renames outerFlange to FTOuterFlange in ten applicable source and runtime CAD
variations and ten GEMC2 SQLite placement rows. The stainless-steel material, color, mother, and passive
placement are retained, including the original 1272.77 mm z translation. That translation differs by 0.5 mm
from the inner-shield and outer-mount placements and is not rounded to their common value.

The original is a closed manifold annular plate with 2,132 triangles and genus 13. The authorized hole
removal closes twelve bolt tunnels with pygemc's remove\_holes utility, leaving the central opening intact
and reducing genus to 1. Flat annular faces are then retriangulated using every original inner and outer rim
coordinate on both z faces. The inner radius is 95.25 mm, the outer radius 127 mm, and the original
thickness approximately 19.049927 mm. Each face retains 97 inner rim vertices and 112 outer rim vertices.
Saved coordinates are verified exactly against those engineering vertices; the original bounds are
unchanged.

Triangles fall to 836 (60.79\% fewer), and bytes fall from 106,684 to 41,884. The output remains one closed
manifold surface with no detected self-intersections or zero-area faces. Closing the bolt holes changes
volume from 389837.402125 to 422180.789284 mm$^3$, an increase of 8.296635\%. Retriangulation changes the
hole-filled reference volume only at floating-point precision. Bidirectional vertex, edge, and surface
sampling against that reference gives a maximum difference of 0.0000148 mm. This sampled difference
describes retriangulation, not the deliberately closed holes, and is not a proven global error bound. The
archived engineering mesh reproduces the optimized STL byte for byte.

Table~\ref{tab:outer-flange-mesh} summarizes the original and hole-filled flange.

\begin{table}[htbp]
\centering
\caption{Original and optimized FT outer flange.}
\label{tab:outer-flange-mesh}
\begin{tabular}{lrr}
\toprule
Quantity & Original & Optimized \\
\midrule
Triangles & 2,132 & 836 \\
Bytes & 106,684 & 41,884 \\
Genus & 13 & 1 \\
Volume (mm$^3$) & 389837.402 & 422180.789 \\
\bottomrule
\end{tabular}
\end{table}

The isolated tracking fixture contains only this stainless-steel flange in air. Both cases use the new name
and identical material and placement attributes. Matched 5 GeV LUND primaries originate at the target,
uniform in solid angle over 1.9--2.8 degrees and full azimuth. FTFP\_BERT runs without optical processes or
magnetic field, with one thread, default production cuts, and no GUI or event output. Five alternating-order
repetitions use seeds 12345--12349 and 200 events for each electron, proton, positive-pion, and neutron run.
Separate zero-event runs measure startup; elapsed time since beamOn measures transport. A fixed executable
snapshot is used for every run, with its hash checked afterward. Earlier runs interrupted by sandbox
resource-timer restrictions or invalidated by a changing shared executable are excluded.

Table~\ref{tab:outer-flange-timing} records the tracking measurements.

\begin{table}[htbp]
\centering
\caption{Median transport wall seconds per 200 events over five repetitions.}
\label{tab:outer-flange-timing}
\begin{tabular}{lrrr}
\toprule
Particle & Original & Optimized & Original/optimized \\
\midrule
e- & 0.646804 & 0.668927 & 0.967 \\
proton & 0.429432 & 0.425280 & 1.010 \\
pi+ & 0.393231 & 0.377190 & 1.043 \\
neutron & 0.221017 & 0.274356 & 0.806 \\
\bottomrule
\end{tabular}
\end{table}

Median startup simulation wall time changes from 0.699890 to 0.693861 s; peak startup RSS changes from
117.83 to 115.67 MiB. Median original/optimized transport ratios are 0.967, 1.010, 1.043, and 0.806 for
electrons, protons, positive pions, and neutrons, respectively. A ratio above one means shorter optimized
tracking time; a ratio below one means longer time. All 52 processes complete normally, and all forty
transport runs report their 200 requested events. These short isolated workloads do not establish a
universal tracking gain. The added stainless steel changes scattering and stochastic shower histories, so
timing differences cannot be attributed solely to triangle reduction or treated as detector-response
equivalence.

Table~\ref{tab:outer-flange-warnings} records isolated-fixture exceptions across all 52 processes.

\begin{table}[htbp]
\centering
\caption{Geant4 exception counts in the timed isolated flange fixtures.}
\label{tab:outer-flange-warnings}
\begin{tabular}{lrr}
\toprule
Diagnostic & Original & Optimized \\
\midrule
GeomSolids1001 & 0 & 0 \\
GeomNav1002 & 0 & 0 \\
GeomNav0003 & 0 & 0 \\
GeomVol1002 & 0 & 0 \\
\bottomrule
\end{tabular}
\end{table}

A separate twelve-volume assembly fixture includes the flange, unchanged inner shield and outer mount, and
nine copper bolts. Both original and optimized fixtures use the new flange name and the same placements.
Their 10,000-point Geant4 checks report 5 and 5 overlap warnings, respectively; both also retain one
surface-hole import warning from the unchanged outer mount. The retained logs identify the volumes and
micrometre-scale penetrations. Existing outer-mount/flange interference remains a limitation of the
assembly, so no claim of complete overlap freedom is made. Full default construction completes without
missing meshes, with 4 solid warnings and 9 overlap reports. All fifteen native geometry/material variations
and fourteen CAD variations pass the semantic comparison against GEMC2. The earlier six beamline studies
retain their independent timing tables, diagnostic results, and archives.

The subsequent mount study resolves this earlier joint interference; see Section~\ref{sec:outer-mount}.

\subsection{Tracking at fixed material volume}

A follow-up control compares the hole-filled flange before and after retriangulation, removing the added
material as a difference between the cases. The baseline has 2,084 triangles and 104,284 bytes; the
simplified flange has 836 triangles and 41,884 bytes, a 59.88\% triangle reduction relative to this filled
reference. Both have the same central opening, bounds, stainless-steel material, placement, and volume of
422180.789284 cubic millimetres to the reported precision. Their sampled surface differences remain below
0.0000148 mm. The fixture label original now denotes the unreduced hole-filled reference, not the drilled
engineering mesh.

The control uses the same fixed executable snapshot and physics settings as the earlier flange study, with 5
GeV primaries uniform in solid angle over 1.9--2.8 degrees and full azimuth. Five alternating-order
repetitions use seeds 12345--12349. Each particle run has 2,000 events, ten times the earlier sample; both
cases have identical primary hashes and geometry/material database records. Separate startup samples retain
the loading and peak-memory measurements. The larger control is reported independently, without pooling its
absolute times or ratios with the earlier 200-event drilled-versus-filled results.

Table~\ref{tab:outer-flange-control} gives median transport wall seconds per 2,000 events.

\begin{table}[htbp]
\centering
\caption{Hole-filled reference versus simplified flange: medians over five repetitions.}
\label{tab:outer-flange-control}
\begin{tabular}{lrrr}
\toprule
Particle & Filled reference & Simplified & Reference/simplified \\
\midrule
e- & 7.40750 & 6.90088 & 1.073 \\
proton & 4.63719 & 4.56330 & 1.016 \\
pi+ & 4.19454 & 4.05502 & 1.034 \\
neutron & 2.94839 & 2.91016 & 1.013 \\
\bottomrule
\end{tabular}
\end{table}

Median baseline/simplified tracking wall-time ratios are 1.073 for electrons, 1.016 for protons, 1.034 for
positive pions, and 1.013 for neutrons. A value above one means shorter tracking time for the simplified
mesh. These measurements quantify retriangulation at fixed material volume. The simplified mesh is faster in
all twenty particle/seed pairs; the individual paired ratios are retained in the control dataset. Matched
primaries and seeds do not guarantee identical subsequent shower histories because tessellated boundary
arithmetic can alter stochastic transport. Five repetitions do not establish a universal speedup or a
precise causal decomposition of the earlier slowdown.

Median startup simulation wall time changes from 0.704797 to 0.692080 s, with peak startup RSS changing from
117.52 to 115.58 MiB. All 52 control processes complete normally, and all forty transport runs report the
requested 2,000 events. Geant4 exception counts, paired seed ratios, fixture hashes, and raw timing reports
are retained in a separate control dataset. Neither case records a Geant4 exception in the sampled control.

\FloatBarrier
Figure~\ref{fig:outer-flange-stl} compares the original and filled triangulations.

\begin{figure}[htbp]
\centering
\includegraphics[width=\textwidth]{figures/outer_flange_stl_comparison.pdf}
\caption{Original FT outer flange (left) and optimized flange with twelve bolt holes filled (right).
Matched cameras and bounds use original STL coordinates; both retain the central opening and rim vertices.}
\label{fig:outer-flange-stl}
\end{figure}
